\documentclass[11pt]{article}
\usepackage[T1]{fontenc}
\usepackage{lmodern,amsmath,amssymb,amsthm,booktabs,array,longtable}
\usepackage[margin=1in]{geometry}
\usepackage[colorlinks=true,linkcolor=blue,citecolor=blue,urlcolor=blue]{hyperref}
\usepackage{xurl}
\newtheorem{theorem}{Theorem}
\newtheorem{lemma}[theorem]{Lemma}
\newtheorem{corollary}[theorem]{Corollary}
\title{Proof of Chen and Chen's parity-positivity conjectures:\\
OEIS A392664--A392666, COMMENTS, and Conjecture 2 (p.~1)}
\author{Haobo Ma\\\small ChronoAI Pte Ltd, Singapore
\and Wenlin Zhang\\\small National University of Singapore\\
\small ChronoAI Pte Ltd, Singapore}
\date{3 November 2026}
\begin{document}
\maketitle
\begin{abstract}
We prove Conjecture~2 on page~1 of Chen and Chen's January 2026
executive summary and the parity-positivity conjectures in OEIS
A392664--A392666, COMMENTS: every positive even row of the step-$m$
Eulerian triangle is strictly positive for real $m\leq-3$.
We also prove the adjacent odd-row alternating-sum nonvanishing
assertions at $m=-4,-5$. For the largest root modulus $M_n(m)$,
we obtain $M_n(-3)\sim5^{n/2}/\sqrt{45}$, correcting the stated
growth exponent $\log2$ to $\tfrac12\log5$ and the proposed scaling
constant $1/3$. The constants $1/6$ and $1/10$ at steps $-4,-5$
are proved. Written arguments also give the integer-step sign
classification, companion identities and boundary cases.
\end{abstract}

\section{The conjectures and the recurrence}
The COMMENTS sections of OEIS A392664--A392666 conjecture
even-row nonnegativity at steps $-3,-4,-5$ \cite{oeis3,oeis4,oeis5}.
Their executive summary states the conjecture for every integer $m\leq-3$
\cite[p.~1, Conjecture~2]{chen} and reports verification through row 400
for the three listed steps. We prove that stronger parameter-wide statement,
with strict positivity in place of nonnegativity.

For a real parameter $m$, define $T_m(n,k)$ for integers $1\leq k\leq n$ by
\begin{align}
 T_m(n,1)&=T_m(n,n)=1 \qquad(n\geq1),\label{eq:boundary}\\
 T_m(n,k)&=\bigl(m(n-k)+1\bigr)T_m(n-1,k-1)
          +\bigl(m(k-1)+1\bigr)T_m(n-1,k),\label{eq:recurrence}
\end{align}
where the second equation applies only for $1<k<n$.
These equations give existence and uniqueness on the triangular domain
by induction on $n$. If $m$ is an integer, all entries are integers.
No values outside the triangle enter the equations.
Set $R_{m,n}(y)=\sum_{k=1}^n T_m(n,k)y^{k-1}$.

\begin{theorem}\label{thm:main}
For every real $r\geq3$, the triangle of step $m=-r$ satisfies
\begin{enumerate}
\item $T_{-r}(n,k)>0$ for every positive even $n$ and $1\leq k\leq n$;
\item $T_{-r}(n,k)<0$ for every odd $n\geq3$ and $1<k<n$.
\end{enumerate}
In particular, the three OEIS conjectures and Conjecture~2 of
\cite{chen} hold for their stated integer parameters.
\end{theorem}

\section{Two-row sign propagation}
First observe that for every real $m$,
\begin{equation}\label{eq:symmetry}
 T_m(n,k)=T_m(n,n+1-k).
\end{equation}
Indeed, reflection interchanges the boundary entries. For an interior
entry, the two predecessor indices $k-1,k$ are reflected in row $n-1$
to $n+1-k,n-k$, respectively. The two coefficients in
\eqref{eq:recurrence} are interchanged as well. Induction on $n$
therefore proves \eqref{eq:symmetry}.

\begin{lemma}\label{lem:second}
For $r\geq3$, $T_{-r}(n,2)>0$ whenever $n\geq2$ is even.
\end{lemma}
\begin{proof}
Write $a_n=T_{-r}(n,2)$ for $n\geq2$. The boundary gives $a_2=1$;
substitution at $k=2$ gives, for $n\geq2$,
\[
 a_{n+1}=1-r(n-1)+(1-r)a_n.
\]
Substituting this identity once more and collecting terms yields
\begin{equation}\label{eq:twostep}
 a_{n+2}=(r-1)^2a_n+r(r-2)n+2-r^2.
\end{equation}
The first step must be treated separately:
\[
 a_4=(r-1)^2+2r(r-2)+2-r^2
     =2r(r-3)+3\geq3.
\]
For $n\geq4$, the remaining term in \eqref{eq:twostep} satisfies
\[
 r(r-2)n+2-r^2
 \geq3r^2-8r+2
 =3(r-3)^2+10(r-3)+5>0.
\]
Thus \eqref{eq:twostep}, starting at $a_4>0$, proves positivity for
all later even $n$. Together with $a_2=1$, this proves the lemma.
\end{proof}

\begin{proof}[Proof of Theorem~\ref{thm:main}]
The even row $n=2$ is $(1,1)$. Suppose an even row $n\geq2$
is strictly positive. For $1<k<n+1$, both coefficients defining
$T_{-r}(n+1,k)$ are strictly negative: their positive integer factors
$n+1-k$ and $k-1$ are at least one, so each coefficient is at most
$1-r<0$. Both predecessor entries lie in row $n$ and are positive.
Every interior entry of row $n+1$ is therefore negative.

In row $n+2$, the boundary entries are 1. The entries at $k=2$
and $k=n+1$ are positive by Lemma~\ref{lem:second} and symmetry.
For $3\leq k\leq n$, both predecessor indices $k-1,k$ are
strictly interior to row $n+1$, so both predecessor entries are
negative. Both recurrence coefficients are again negative. Their
products, and hence their sum, are positive. This last index range
is empty when $n=2$, in which case the boundary and second-column
arguments cover the entire row. Induction proves positivity of
all even rows, and the intervening argument proves negativity of
all odd-row interior entries.
\end{proof}

\begin{corollary}\label{cor:classification}
For integer $m$, every positive even row is entrywise nonnegative
if and only if $m\leq-3$ or $m\geq0$. For exactly these parameters,
every positive even row is in fact strictly positive.
\end{corollary}
\begin{proof}
Theorem~\ref{thm:main} covers $m\leq-3$. If $m\geq0$, both
coefficients in \eqref{eq:recurrence} are at least one, so induction
gives strict positivity of every row. The two remaining integer
parameters fail already at
\[
 T_m(4,2)=T_m(4,3)=2m^2+6m+3=-1
 \qquad(m=-1,-2).
\]
\end{proof}

\section{Companion identities and root bounds}
These consequences address the adjacent COMMENTS statements in the three
OEIS entries. They concern the same recurrence, rather than a different
root-asymptotic assertion.

\begin{corollary}\label{cor:sums}
For every real $m$,
\[
 S_n(m):=\sum_{k=1}^n T_m(n,k)
 =\prod_{j=0}^{n-2}(mj+2),\qquad S_1(m)=1.
\]
For even $n$,
\[
 R_{m,n}(-1)=0,\qquad
 \sum_{\substack{1\leq k\leq n\\ k\text{ odd}}}T_m(n,k)
 =\sum_{\substack{1\leq k\leq n\\ k\text{ even}}}T_m(n,k)
 =\tfrac12 S_n(m).
\]
\end{corollary}
\begin{proof}
For $n\geq3$, extend the preceding row by zeros just for summing
\eqref{eq:recurrence}. At $k=1,n$, the same formula then gives
the boundary value 1. A preceding-row entry at position $j$
has total coefficient
\[
 \bigl(m(n-j-1)+1\bigr)+\bigl(m(j-1)+1\bigr)=m(n-2)+2.
\]
Consequently $S_n=(m(n-2)+2)S_{n-1}$. Since $S_2=2$,
the product follows. For even $n$, reflection pairs equal entries
at positions $k$ and $n+1-k$ of opposite parity. Their contributions
to $R_{m,n}(-1)$ cancel, and their parity sums agree.
\end{proof}

\begin{corollary}\label{cor:annulus}
For real $m\leq-3$ and even $n\geq2$, every root $z$ of
$R_{m,n}$ satisfies
\[
 L_n\leq |z|\leq U_n,\qquad
 L_n=\min_{1\leq k<n}\frac{T_m(n,k)}{T_m(n,k+1)},\quad
 U_n=\max_{1\leq k<n}\frac{T_m(n,k)}{T_m(n,k+1)}.
\]
All these ratios are positive; symmetry also gives $L_nU_n=1$.
\end{corollary}
\begin{proof}
This is the Enestr\"om--Kakeya bound \cite{kakeya}; the following
argument includes the boundary cases. For a polynomial
$P(z)=\sum_{j=0}^d a_jz^j$ with $a_j>0$, put
$U=\max_{0\leq j<d}a_j/a_{j+1}$ and $b_j=a_jU^j$.
Then $b_0\leq\cdots\leq b_d$. If
$Q(w)=P(Uw)$ had a root with $|w|>1$, the identity
\[
 (1-w)Q(w)=b_0+\sum_{j=1}^d(b_j-b_{j-1})w^j-b_dw^{d+1}
\]
would imply
\[
 b_d|w|^{d+1}
 \leq b_0+\sum_{j=1}^d(b_j-b_{j-1})|w|^j
 \leq b_d|w|^d,
\]
a contradiction. Applying this upper bound to $z^dP(1/z)$
gives the lower bound. Theorem~\ref{thm:main} supplies positive
coefficients here. Reflection makes the set of consecutive ratios
invariant under reciprocation, so $L_n=1/U_n$.
\end{proof}

\section{Odd-row alternating sums}
Put $A_q(m)=R_{m,2q+1}(-1)$, with $A_0(m)=1$.
The COMMENTS of A392665 and A392666 assert nonvanishing for
every odd row, not just the displayed examples.

\begin{corollary}\label{cor:odd}
For every integer $q\geq1$,
\[
 A_q(-4)/8^q\ \hbox{is an odd integer},\qquad
 A_q(-5)\equiv1\pmod3.
\]
More generally, the first conclusion holds with $-4,8$ replaced by
$-r,2r$ for every positive integer $r$ divisible by $4$; the second
holds for every integer $m\equiv1\pmod3$. In particular, both
source nonvanishing assertions hold for all odd rows.
\end{corollary}
\begin{proof}
Multiplication of the triangle recurrence by $y^{k-1}$ and summation gives
\begin{equation}\label{eq:poly}
 R_{m,n+1}(y)=[1+(m(n-1)+1)y]R_{m,n}(y)
              +my(1-y)R'_{m,n}(y).
\end{equation}
For $m\ne0$ the exponential generating function is
\begin{equation}\label{eq:egf}
 H(t,y):=\sum_{j\geq0}R_{m,j+1}(y)\frac{t^j}{j!}
 =e^{(1-y)t}\left(\frac{1-y}{1-ye^{m(1-y)t}}\right)^{2/m}.
\end{equation}
Indeed, differentiation verifies
$(1-myt)H_t=(1+y)H+my(1-y)H_y$ and $H(0,y)=1$.
Coefficient comparison gives \eqref{eq:poly}, so uniqueness proves
the formula. It is a Taylor series at $t=0$ for $y\ne1$.
At $y=-1$ it gives
\[
 \sum_{q\geq0}A_q(-r)\frac{t^{2q}}{(2q)!}
       =\cosh(rt)^{2/r}.
\]
Set $C_j=(d/ds)^{2j}\log\cosh(s)|_{s=0}$.
These are integers and $C_1=1$: successive derivatives of
$\tanh s$ are integer polynomials in $\tanh s$, by
$(\tanh s)'=1-\tanh^2s$.
The logarithm of the last generating function has even derivatives
$\kappa_{2j}=2r^{2j-1}C_j$ and odd derivatives zero.
Expanding its exponential expresses $A_q(-r)$ as the sum over
partitions of $2q$ labelled elements into even blocks, with weight
$\prod_B\kappa_{|B|}$. After division by $(2r)^q$, a block of
size $2j$ contributes $(r/2)^{j-1}C_j$.
When $4\mid r$, every partition with a block larger than two
therefore contributes an even integer. The remaining
$(2q-1)!!$ pair partitions each contribute 1. This proves the
first assertion.

If $m\equiv1\pmod3$, reduction of the recurrence gives
$R_{m,n}\equiv R_{1,n}\pmod3$. The Eulerian identity
\[
 R_{1,n}(y)=(1-y)^{n+1}\sum_{j\geq0}(j+1)^n y^j
\]
follows by applying $(y\,d/dy)^n$ to $y/(1-y)$ and comparing
\eqref{eq:poly}. For odd $n$, $(j+1)^n\equiv j+1\pmod3$,
so this formal identity gives
$R_{m,n}(y)\equiv(1-y)^{n-1}\pmod3$.
At $y=-1$ the value is $2^{n-1}\equiv1\pmod3$, as required.
\end{proof}

\section{Root growth and the source's limiting constants}
For $n\geq2$, let $M_n(m)$ be the largest modulus among the roots
of $R_{m,n}$. The following corrects the $m=-3$ instance of
\cite[p.~1, Theorem~1; pp.~2--3, tables]{chen}, and proves the
two other listed negative-step limiting constants.

\begin{theorem}\label{thm:growth}
For every real $m>0$ and every real $r\geq4$,
\[
 \frac{M_n(m)}{(m+1)^n}\longrightarrow\frac{2}{m(m+1)},\qquad
 \frac{M_n(-r)}{(r-1)^n}\longrightarrow\frac{2}{r(r-1)}.
\]
At step $-3$, instead,
\[
 \frac{M_n(-3)}{5^{n/2}}\longrightarrow\frac1{\sqrt{45}},\qquad
 \lim_{n\to\infty}\frac{\log M_n(-3)}n=\tfrac12\log5.
\]
In particular, $M_n(-4)/3^n\to1/6$ and
$M_n(-5)/4^n\to1/10$, whereas $M_n(-3)/2^n\to\infty$.
The proposed growth invariance between steps $1$ and $-3$ fails.
\end{theorem}
\begin{proof}
Write $a_j(n)=T_m(n,j+1)$; reflection also makes it the
coefficient of $y^{n-1-j}$. Solving the second-column recurrence
gives, for $m\ne0$ and $n\geq2$,
\begin{equation}\label{eq:column2}
 a_1(n)=\frac{2(m+1)^{n-1}}m-n+1-\frac2m.
\end{equation}
Substitution and $a_1(2)=1$ verify it. At $m=-3$, the next column
satisfies $a_2(n+1)=(7-3n)a_1(n)-5a_2(n)$, with $a_2(3)=1$.
Consequently
\begin{equation}\label{eq:column3}
 a_2(n)=\frac{(-5)^n}{45}
       +\left(\frac59-\frac n3\right)(-2)^n
       +\frac{n^2}2-\frac{13n}6+\frac{20}9
       \quad(n\geq3),
\end{equation}
again verified by substitution and the initial value.

Here is a uniform bound for the remaining columns. Extend each row
by zeros outside its triangular domain. The same recurrence then
holds at its endpoints. A path from row 1 to position $j+1$ of
row $n$ has $j$ index-increasing steps and $n-1-j$ other steps.
There are $\binom{n-1}{j}$ such paths. Every increasing-step weight
has modulus at most $(|m|+1)n$. The other weights have modulus at
most $B_j=1+mj$ for $m>0$, or $B_j=rj-1$ for $m=-r$, $r\geq3$.
Therefore, for $1\leq j\leq n-1$,
\begin{equation}\label{eq:pathbound}
 |a_j(n)|\leq[(|m|+1)n^2]^j B_j^{\,n-1-j}.
\end{equation}
For $m>0$, $B_j^{1/j}\leq\sqrt{1+2m}$ for $j\geq2$:
$\log(1+mx)/x$ decreases for $x>0$, since
$\log(1+v)>v/(1+v)$ for $v>0$.
For $r\geq4$, $B_j^{1/j}\leq\sqrt{2r-1}$ for $j\geq2$;
here $\log(rx-1)/x$ decreases for $x\geq2$, since
$\log(rx-1)>rx/(rx-1)$ when $rx-1\geq7$.
For $r=3$ and $j\geq3$, $3j-1\leq2^j$ by induction.

In the first two cases put respectively $L_n=(m+1)^n$ and
$L_n=(r-1)^n$. Since
$\sqrt{1+2m}<m+1$ and $\sqrt{2r-1}<r-1$ for these parameters,
\eqref{eq:pathbound} implies
\begin{equation}\label{eq:tail}
 |a_j(n)|/L_n^j\leq\epsilon_n^j\quad(j\geq2),
 \qquad \epsilon_n\longrightarrow0.
\end{equation}
One may take $\epsilon_n=(|m|+1)n^2(b/\lambda)^n$,
where $(b,\lambda)$ is the corresponding pair of bases above.
On every region $|z|\geq\delta>0$, the equation for a nonzero
root $y=L_nz$, divided by $L_n^{n-1}z^{n-2}$, is
\[
 z+\frac{a_1(n)}{L_n}
    +\sum_{j=2}^{n-1}\frac{a_j(n)}{L_n^j}z^{1-j}=0.
\]
The last sum is bounded by
$\delta\sum_{j\geq2}(\epsilon_n/\delta)^j\to0$, uniformly on
that region. By \eqref{eq:column2}, the linear part tends to
$z+c$, with $c=2/[m(m+1)]$ for $m>0$, or to
$z+(-1)^nc$, with $c=2/[r(r-1)]$ for $m=-r$.
On a small circle about its nonzero limiting zero, the error is
smaller than the linear part for large $n$.
Rouch\'e's theorem \cite[\S1.10(iv)]{dlmf} gives a root there.
For $|z|\geq c+\eta$, the linear part has modulus at least
$\eta/2$ for large $n$, and the same estimate excludes roots.
Thus the largest scaled root modulus tends to $c$.

For $m=-3$ put $L_n=5^{n/2}$. The final column bound above gives
\eqref{eq:tail} for $j\geq3$, with
$\epsilon_n=4n^2(2/\sqrt5)^n$.
Divide the root equation instead by $L_n^{n-1}z^{n-3}$.
The same geometric estimate makes its tail tend uniformly to zero
on $|z|\geq\delta$. Equations \eqref{eq:column2}--\eqref{eq:column3}
give the limiting quadratic
\[
 z^2+(-1)^n/45.
\]
Its two simple zeros have modulus $1/\sqrt{45}$. Apply Rouch\'e's
theorem on small circles about them, separately for even and odd
$n$. For $|z|\geq1/\sqrt{45}+\eta$, the full quadratic part
has modulus at least
\[
 |z|^2-|a_1(n)/L_n|\,|z|-|a_2(n)/L_n^2|,
\]
which is bounded away from zero for large $n$.
The uniform tail estimate therefore excludes exterior roots.
This proves the remaining limit and its logarithmic consequence.
At $m=1$ the first limit gives exponent $\log2$, so the stated
mirror invariance is false for this pair.
\end{proof}

For a concrete finite bound, the accompanying exact checker gives
\[
 5^{200}/7<M_{400}(-3)<5^{200}/6.
\]
On $|z|=1/7$ and $1/6$, the quadratic $z^2+1/45$ has modulus
at least $4/2205$ and $1/180$, respectively, whereas the preceding
error bound is less than $10^{-6}$. Comparing with
$z^{397}(z^2+1/45)$ gives 397 zeros inside the first circle and
all 399 inside the second. Exactly two lie between them.

The $m=-1$ coefficient formula already appears in Greubel's
1 March 2022 addition to OEIS A144431, FORMULA \cite{oeislinked}.
Its factorization gives the following boundary-case root description.
\begin{corollary}\label{cor:boundaryroots}
For $q\geq1$, the two excluded integer steps satisfy
\[
 R_{-1,2q}(y)=R_{-2,2q}(y)=(1-y)^{2q-2}(1+y),
\]
and
\[
 R_{-1,2q+1}(y)=(1-y)^{2q-2}(1+y^2),\qquad
 R_{-2,2q+1}(y)=(1-y)^{2q}.
\]
Thus $M_n(-1)=M_n(-2)=1$ for $n\geq2$; both growth exponents
are zero. Also $R_{0,n}=(1+y)^{n-1}$. For $m>0$ and $n\geq2$,
every root of $R_{m,n}$ is simple and negative. At $m=-3$,
nonreal roots occur already in row 6.
\end{corollary}
\begin{proof}
The factor formulas follow from \eqref{eq:poly} in successive
even--odd pairs, starting with $R_{m,2}=1+y$ and
$R_{-1,3}=1+y^2$, $R_{-2,3}=(1-y)^2$.
At $m=0$, \eqref{eq:poly} multiplies by $1+y$ at each step.
For $m>0$, at a negative simple zero $x$ of $R_{m,n}$,
\eqref{eq:poly} has sign opposite to $R'_{m,n}(x)$.
The intermediate value theorem gives one zero between each pair
of consecutive negative zeros, one between the rightmost zero
and 0, and one to the left of the leftmost zero (the leading
coefficient is 1). These $n$ distinct zeros exhaust degree $n$;
induction starts at $R_{m,2}=1+y$.
Finally,
\[
 R_{-3,6}(y)=(y+1)(y^4+16y^3+246y^2+16y+1).
\]
For real $y\ne0$, the quartic divided by $y^2$ is
$(y+y^{-1}+8)^2+180>0$, so its four roots are nonreal.
\end{proof}
The zero exponent at $m=-1$ corrects the summary's
page-2 entry $\beta=-\infty$. A monic polynomial with constant
term 1 has product of root moduli 1, hence $M_n\geq1$.
This also rules out the source's formula $\beta(m)=\log|m+1|$
when $0<|m+1|<1$. For $|m+1|>1$, the root $-1$ in every
even row scales to modulus $|m+1|^{-n}\to0$: an annulus with a
fixed positive inner radius cannot contain \emph{all} scaled zeros.
The finite-row annulus of Corollary~\ref{cor:annulus} remains valid.

\section{Coverage and verification}
\small
\begin{longtable}{@{}p{.46\linewidth}p{.47\linewidth}@{}}
\toprule
Source item & Coverage \\
\midrule
\endhead
A392664--A392666, COMMENTS: even-row conjectures;
\cite[p.~1, Conjecture~2]{chen}
& Proved: Theorem~\ref{thm:main}, including all real $m\leq-3$. \\
Same entries: reflection, row sums, even-row alternating sums
& Proved: \eqref{eq:symmetry} and Corollary~\ref{cor:sums}. \\
Same entries: coefficient-ratio root annulus
& Proved: classical bound cited and proved in Corollary~\ref{cor:annulus}. \\
Cross-referenced steps $m=-2,-1,0,1,2,3$
& Proved: Corollaries~\ref{cor:classification}, \ref{cor:boundaryroots}
and Theorem~\ref{thm:growth}; recurrence identifications appear
in the source guide. \\
A392665, A392666, COMMENTS: odd-row nonvanishing
& Proved: Corollary~\ref{cor:odd}, with stronger congruences. \\
A392664, COMMENTS: $R_{-3,5}(-1)=0$
& Proved by \eqref{eq:egf} or row 5; displayed odd-row values
are also exact finite checks. \\
\cite[pp.~2--3]{chen}: constants $1/6,1/10$ at $m=-4,-5$
& Proved: Theorem~\ref{thm:growth} covers every real $r\geq4$. \\
\cite[p.~1, Theorem~1; pp.~2--3]{chen}: step $-3$ growth,
constant $1/3$, and growth invariance with $m=1$
& Refuted-and-corrected: Theorem~\ref{thm:growth} gives exponent
$\tfrac12\log5$ and scale $5^{n/2}/\sqrt{45}$. \\
\cite[p.~1, Theorem~1; p.~2]{chen}: unrestricted complex growth
law and $m=-1$ degeneracy
& Refuted-and-corrected: $M_n\geq1$; the exponent at $m=-1$ is 0.
Remaining complex-parameter asymptotics are not covered. \\
\cite[pp.~2--3]{chen}: nonreal roots at $m=-3$; all roots real
at $m=-4,-5$
& Proved at $m=-3$ by row 6; all-row reality at $m=-4,-5$
is not covered. \\
\cite[p.~2]{chen}: fixed annulus for scaled zeros
& Refuted-and-corrected as a uniform positive inner radius for
all zeros; the finite-row bound above applies. \\
OEIS mirror map $m\mapsto-m-2$ and equality of $|m+1|$
& Proved by substitution; it does not imply equal growth
exponents, as steps $1,-3$ show. \\
\cite[pp.~2--3]{chen}: finite plots and EK tightness percentages;
page-3 reference to a full preprint
& Not covered as numerical accuracy claims or as unprovided
additional theorems; no conclusion depends on them. \\
\bottomrule
\end{longtable}
\normalsize

The pinned
\href{https://github.com/the-omega-institute/trureturing/blob/66a0495cbdf1bb9b42d715c1b3cd3365862896fa/D5/S1/Recurrence/Algebraic/ChenChenParityPositivity.lean}{\texttt{ChenChenParityPositivity.lean}}
at its 2 November 2026 revision contains the declarations
\texttt{even\_rows\_positive} and \texttt{result}, in namespace
\nolinkurl{D5.S1.Recurrence.Algebraic.ChenChenParityPositivity}.
The first formalizes strict positivity for integer $r\geq3$ and
every function satisfying the boundary and recurrence. The second
formalizes existence together with universal even-row nonnegativity
for steps $-3,-4,-5$. The real-parameter extension, integer-step
classification, identities, alternating-sum congruences and root-growth
results have written proofs above; they are not additional assertions
of either pinned declaration.
The accompanying \texttt{verify.py} checks exact integer rows through
$n=400$ for $m=-12,\ldots,-3,0,1,2,3$: symmetry, the row-sum
product, even-row positivity and parity sums, and negative-step
odd-row interior signs. It also checks the displayed initial values
and the two excluded integer parameters. Further exact checks cover
the alternating-sum congruences, column formulas and degenerate-step
factorizations. These finite computations are separate from the all-row proofs.

\clearpage
\begin{thebibliography}{9}
\bibitem{oeislinked}
The On-Line Encyclopedia of Integer Sequences,
A225372, A144431, A007318, A008292, A060187 and A142458,
NAME and FORMULA; G.~C.~Greubel's A144431 formula,
1 March 2022.
\url{https://oeis.org/A144431}; links to all six entries
are given in the accompanying source guide.
\bibitem{dlmf}
NIST Digital Library of Mathematical Functions,
\S1.10(iv), Rouch\'e's theorem.
\url{https://dlmf.nist.gov/1.10#iv}.
\bibitem{oeis3}
X.~Chen and Y.~Chen, OEIS A392664, COMMENTS and FORMULA,
18 January 2026; entry revision 35, 5 February 2026.
\url{https://oeis.org/A392664}.
\bibitem{oeis4}
X.~Chen and Y.~Chen, OEIS A392665, COMMENTS and FORMULA,
18 January 2026; entry revision 20, 5 February 2026.
\url{https://oeis.org/A392665}.
\bibitem{oeis5}
X.~Chen and Y.~Chen, OEIS A392666, COMMENTS and FORMULA,
1 February 2026; entry revision 17, 5 February 2026.
\url{https://oeis.org/A392666}.
\bibitem{chen}
X.~Chen and Y.~Chen,
\emph{The Step-$m$ Eulerian Triangle: Mirror Symmetry and Universal
Growth Rate---Executive Summary for OEIS}, January 2026,
p.~1, Conjecture~2. Zenodo version 1, 31 January 2026,
\url{https://doi.org/10.5281/zenodo.18446568}; version 2,
1 February 2026, \url{https://doi.org/10.5281/zenodo.18452218}
(same summary, additional data).
\bibitem{kakeya}
S.~Kakeya, On the limits of the roots of an algebraic equation with
positive coefficients, \emph{T\^ohoku Mathematical Journal}
\textbf{2} (1912), 140--142.
\url{https://www.jstage.jst.go.jp/article/tmj1911/2/0/2_0_140/_article}.
\end{thebibliography}
\medskip\noindent\small
\copyright\ 2026 Haobo Ma and Wenlin Zhang.
Licensed under Creative Commons Attribution 4.0 International (CC BY 4.0).
\end{document}
